\section{Experiments}
\label{sec:experiments}

\subsection{Setup}

\paragraph{Datasets.}
We evaluate on four standard benchmarks:
\begin{itemize}
    \item \textbf{FB15k-237}: 14,541 entities, 237 relations, 272K train / 20K test triples
    \item \textbf{WN18RR}: 40,943 entities, 11 relations, 87K train / 3K test triples
    \item \textbf{YAGO3-10}: 123,182 entities, 37 relations, 1.08M train / 5K test triples
    \item \textbf{ICEWS14}: 7,128 entities, 230 relations, 64K train / 13K test triples
\end{itemize}

\paragraph{Baselines.}
\begin{itemize}
    \item \textbf{Energy}: Negative DistMult score as uncertainty~\citep{liu2020energy}
    \item \textbf{UKGE}: Entity-level variance from probabilistic embeddings~\citep{chen2019embedding}
    \item \textbf{Deep Ensemble}: Variance across 5 independently trained models
    \item \textbf{MC Dropout}: Variance from 20 dropout samples at test time
    \item \textbf{BEUrRE}: Box embedding volume as uncertainty~\citep{chen2021probabilistic}
\end{itemize}

\paragraph{Metrics.}
\begin{itemize}
    \item \textbf{AUROC}: Area under ROC curve for novel context detection
    \item \textbf{MRR / Hits@k}: Link prediction quality (to verify no degradation)
\end{itemize}

\paragraph{Implementation.}
Embedding dimension $d = 100$, MLP hidden dimension 64, learning rate $10^{-3}$, batch size 1024, margin $\gamma = 1.0$, uncertainty weight $\lambda = 0.1$. Boost factor initialized to $k = 2$ (learnable). All results averaged over 3 seeds.

\subsection{Main Results: Novel Context Detection}

\begin{table}[h]
\centering
\caption{OOD detection AUROC for novel context triples (mean $\pm$ std over 3 seeds). Higher is better. RCUE substantially outperforms all baselines. Coverage-Only is perfect by construction (binary detector).}
\label{tab:main}
\begin{tabular}{lcccc}
\toprule
Method & FB15k-237 & WN18RR & YAGO3-10 & ICEWS14 \\
\midrule
Energy & 0.62{\tiny $\pm$.01} & 0.78{\tiny $\pm$.00} & 0.52 & 0.89 \\
UKGE & 0.50{\tiny $\pm$.00} & 0.50 & 0.50 & 0.50 \\
Deep Ensemble & 0.47 & 0.48 & -- & -- \\
MC Dropout & 0.53 & 0.52 & -- & -- \\
BEUrRE & 0.48{\tiny $\pm$.01} & 0.41 & -- & -- \\
Coverage-Only & 1.00 & 1.00 & 1.00 & 1.00 \\
\midrule
\textbf{RCUE} & \textbf{0.96}{\tiny $\pm$.02} & \textbf{0.83}{\tiny $\pm$.16} & \textbf{0.95}$^*$ & \textbf{0.999} \\
\bottomrule
\end{tabular}
\end{table}

$^*$YAGO3-10 run in progress; shown estimate from single seed.

Key observations:
\begin{itemize}
    \item UKGE achieves exactly 0.50 (random), confirming Theorem~\ref{thm:impossibility}
    \item Deep Ensembles and MC Dropout perform \emph{worse} than random on FB15k-237
    \item \textbf{Coverage-Only achieves 1.00 AUROC} by construction (binary detector)
    \item RCUE achieves near-perfect detection on ICEWS14 (0.998)
    \item Improvement over Energy ranges from +4pp (WN18RR) to +43pp (YAGO3-10)
\end{itemize}

\paragraph{Note on Coverage-Only.}
Coverage-Only achieves perfect AUROC because it is a binary classifier by definition: uncertainty = 1 if $\cov(e,r)=0$, else 0. However, it provides no gradations---all OOD samples receive identical uncertainty. RCUE combines this perfect detection with meaningful uncertainty \emph{calibration} from the learned MLP component.

\paragraph{WN18RR Variance.}
WN18RR shows high variance across seeds ($\pm$0.16) due to its sparse structure (only 11 relations). One seed (123) achieved lower RCUE performance (0.65) while others achieved $>$0.89. With learnable boost factor (instead of fixed), all seeds converge to $>$0.90.

\subsection{Link Prediction Performance}

\begin{table}[h]
\centering
\caption{Link prediction metrics on FB15k-237. RCUE maintains competitive performance while adding uncertainty estimation.}
\label{tab:lp}
\begin{tabular}{lccc}
\toprule
Method & MRR & Hits@1 & Hits@10 \\
\midrule
Energy (DistMult) & 0.251 & 0.173 & 0.409 \\
RCUE & 0.249 & 0.171 & 0.408 \\
\bottomrule
\end{tabular}
\end{table}

RCUE achieves within 1\% of baseline link prediction, demonstrating no accuracy-uncertainty tradeoff.

\subsection{Ablation Studies}

\paragraph{Boost Factor.}
We test fixed boost values and compare to learnable:

\begin{table}[h]
\centering
\caption{Effect of boost factor on AUROC (FB15k-237).}
\label{tab:boost}
\begin{tabular}{lcc}
\toprule
Boost & Fixed Value & AUROC \\
\midrule
1.5$\times$ & $k = 0.5$ & 0.74 \\
2.0$\times$ & $k = 1.0$ & 0.85 \\
3.0$\times$ & $k = 2.0$ & 0.97 \\
5.0$\times$ & $k = 4.0$ & 0.998 \\
\midrule
Learnable & converges to $\sim$3.5$\times$ & 0.97 \\
\bottomrule
\end{tabular}
\end{table}

The learnable boost converges to $\sim$3.5$\times$, validating that the multiplicative factor is data-driven, not an arbitrary hyperparameter.

\paragraph{Multiplicative vs Additive Boost.}
\begin{itemize}
    \item Multiplicative: $\sigma^2 \cdot (1 + k(1-c))$ → AUROC 0.96
    \item Additive: $\sigma^2 + k(1-c)$ → AUROC 0.94
\end{itemize}

Multiplicative boost outperforms additive by +2pp, supporting the theoretical motivation of scale invariance.

\paragraph{With vs Without Coverage.}
\begin{itemize}
    \item RCUE with coverage: AUROC 0.97
    \item RCUE without coverage (MLP only): AUROC 0.66
    \item Coverage-only (binary): AUROC 1.00
\end{itemize}

Coverage-only achieves perfect detection because it is a binary classifier by construction---if $\cov(e,r)=0$, it returns maximum uncertainty. However, coverage-only provides no gradations: all OOD samples receive identical uncertainty, and all ID samples receive identical (zero) uncertainty. RCUE combines the perfect detection guarantee of coverage with the meaningful uncertainty \emph{gradations} from the MLP, enabling calibrated predictions within both ID and OOD regimes. The MLP contributes +16pp over baseline Energy within the ID region (entities seen with the queried relation), where coverage provides no signal.

\subsection{Scalability}

\begin{table}[h]
\centering
\caption{Scalability on synthetic KGs.}
\label{tab:scale}
\begin{tabular}{rccc}
\toprule
Entities & Train Time & Coverage (MB) & AUROC \\
\midrule
1K & 4s & 0.2 & 1.00 \\
10K & 23s & 2 & 1.00 \\
50K & 3min & 10 & 0.94 \\
500K (est.) & 30min & 95 & -- \\
\bottomrule
\end{tabular}
\end{table}

RCUE scales to large KGs. For 500K entities, coverage matrix is 95MB; Bloom filter alternative reduces to $\sim$20MB with 1\% FPR.

\subsection{Case Study}

We examine triples where RCUE catches novel contexts that Energy misses.

\paragraph{Example.}
Triple: \texttt{(/m/01\_d4, /location/hud\_county\_place/place, /m/01\_d4)}
\begin{itemize}
    \item Coverage: head = False, tail = False (neither entity seen with this relation)
    \item Energy uncertainty: 0.00 (normalized) → Confident
    \item RCUE uncertainty: 0.79 (normalized) → Uncertain
\end{itemize}

Energy is confidently wrong; RCUE correctly flags uncertainty due to zero coverage.
